\chapter{Đánh giá và kết luận}

\section{Kết quả đạt được}

\begin{table}[htbp]
\centering
\caption{Tổng hợp kết quả các thành phần}
\label{tab:results}
\begin{tabular}{>{\raggedright\arraybackslash}p{0.25\linewidth} >{\raggedright\arraybackslash}p{0.65\linewidth}}
\toprule
\rowcolor{headerblue} \textcolor{white}{\textbf{Thành phần}} & \textcolor{white}{\textbf{Kết quả}} \\
\midrule
\rowcolor{rowgray} ESP32-CAM & Dual-core: audio 16 kHz liên tục + video JPEG streaming, điều khiển camera real-time (21 tham số) \\
Khử nhiễu DTLN & 1,8M tham số, 3.9 MB, 32 ms latency, huấn luyện trên VIVOS tiếng Việt \\
\rowcolor{rowgray} Voice Rephraze & STT qua gpt-4o-transcribe + sinh mô tả qua GPT-4o-mini, 4+ phong cách viết \\
Phân loại sản phẩm & Accuracy 90,14\%, F1 90,00\%, 32 danh mục, 45 ms/request (CPU) \\
\rowcolor{rowgray} Server trung gian & Relay TCP↔WebSocket, proxy API, spectrogram visualization \\
\bottomrule
\end{tabular}
\end{table}

\section{Ưu điểm}

\begin{itemize}
\item
  \textbf{Pipeline end-to-end:} Từ thu âm ESP32 → khử nhiễu → nhận dạng giọng nói
  → sinh mô tả → phân loại sản phẩm, hoạt động liền mạch.
\item
  \textbf{DTLN nhỏ gọn:} Chỉ 3,9 MB (1,8M tham số), xử lý real-time 32 ms,
  phù hợp thiết bị nhúng tài nguyên hạn chế.
\item
  \textbf{Dual-domain denoising:} Kết hợp STFT (miền tần số) và Conv1D Encoder
  (miền đặc trưng học được) khử nhiễu toàn diện hơn phương pháp đơn miền.
\item
  \textbf{Phân loại chính xác:} XLM-RoBERTa đạt 90,14\% accuracy, hỗ trợ
  đa ngôn ngữ nhờ nền tảng cross-lingual.
\item
  \textbf{Dual-core ESP32:} Tách riêng audio/video trên 2 core đảm bảo
  không gián đoạn stream, 16 DMA buffer cho audio ổn định.
\item
  \textbf{Sinh mô tả linh hoạt:} Nhiều phong cách viết, prompt engineering
  xử lý tốt text lộn xộn từ STT.
\end{itemize}

\section{Hạn chế}

\begin{itemize}
\item
  \textbf{DTLN chỉ mono channel:} Chưa khai thác beamforming hay mảng microphone.
\item
  \textbf{Phụ thuộc cloud cho STT:} Voice Rephraze dùng OpenAI API,
  yêu cầu internet và phát sinh chi phí.
\item
  \textbf{DTLN nhạy cảm nhiễu phi tĩnh:} Hiệu năng giảm với nhiễu
  biến đổi nhanh hoặc âm nhạc phức tạp.
\item
  \textbf{Pha không tối ưu:} Giai đoạn 1 tái sử dụng pha gốc, có thể
  gây artifact khi nhiễu nặng.
\item
  \textbf{ESP32 WiFi dependency:} Thiết bị phụ thuộc kết nối WiFi ổn định,
  không có buffer offline.
\end{itemize}

\section{Kết luận}

Hệ thống ShopNexus đã xây dựng thành công pipeline xử lý AI hoàn chỉnh
cho bài toán hỗ trợ bán hàng bằng giọng nói trên nền tảng nhúng.
Thiết bị ESP32-CAM khai thác hiệu quả kiến trúc dual-core để thu đồng thời
âm thanh và hình ảnh. Mô hình DTLN được huấn luyện trên dữ liệu tiếng Việt
VIVOS cho kết quả khử nhiễu tốt với latency phù hợp real-time.
Dịch vụ phân loại sản phẩm đạt accuracy 90,14\% trên 32 danh mục,
đáp ứng yêu cầu ứng dụng thực tế.

\section{Hướng phát triển}

\begin{itemize}
\item Triển khai PhoWhisper local để giảm phụ thuộc cloud API và chi phí.
\item Bổ sung multi-microphone beamforming cho ESP32.
\item Thêm windowing function trong STFT của DTLN để giảm spectral leakage.
\item Tối ưu mô hình phân loại cho tiếng Việt (fine-tune thêm trên dữ liệu Việt).
\item Triển khai DTLN inference trực tiếp trên Raspberry Pi 4 với TF Lite.
\end{itemize}
